\chapter{调查结果及分析}

% 笔者仔细研读了《课标2017》中对于数学运算素养水平划分的二维结构框架， 在此框架基础之上制定了问卷。 通过对西藏拉萨中学高二年级的12 个班级进行问卷调查。 共发放问卷480份，回收问卷436 份，回收率90.8\%, 其中有效问卷389份，有效问卷占冋收问卷的89.2\%。

% 本章将整理和分析这389份有效问卷。

作者认真研究了《课标2017》所载的关于学生的数学计算能力的2-D结构体系，并以该体系为依据编制了一项调查表.在西藏拉萨高中12个班进行了问卷调查，共计480个学生，436个，回收率90.8\%.

这一章将对389个有效的调查表进行分类和剖析.

\section{问卷信效度分析}

在对调查结果进行分析前，对该问卷进行了信度与效度的检验。如果问卷具有良好的信效度，说明调查是可行的；反之，则不能被接受。因此，需要对问卷进行适当地修改，使其更适合于测量学生的学业成就。 本研究使用SPSS24对调查问卷的信度、效度进行了研究.


\subsection{效度分析}

本文从调查的前期准备到实施调查，采用如下方法来提高调查的有效性。


\begin{enumerate}[(1)]
  \item 前期准备 \par
  作者认真研究了《课标2017》中关于学生的计算能力层次的2-D结构体系，并据此制定了调查研究的基本架构.
  \item 设置题目 \par
  在编写调查问卷的初始阶段，首先按照调查的基本结构，对水平1和水平2的四个方向的四个部分进行三个问题，一共24个问题。
  \item 生成问卷 \par
  为保证问卷调查的正确性，对24道题的次序进行了随机的修改，以避免出现错误。
  \item 发放问卷 \par
  由实习导师预先与教学部门联系，以「校方调查表」的观点，由各班的班长到班监部进行问卷调查，于三月六日的班会上进行，由高二16个班抽取，随机选取12个班，从中选取480人，以保证问卷调查的准确性
  \item 实施过程 \par
  SPSS法用KMO法和巴特利特法进行了比较，结果显示在表~\ref{table:问卷效度分析：KMO和巴特利特检验}中。
  % 通过因子分析检验问卷的有效性。SPSS分析采用KMO和巴特利特检验， 结果见表.
  
  \begin{table}[htbp]
    \centering
    \caption{问卷效度分析：KMO和巴特利特检验}
    \label{table:问卷效度分析：KMO和巴特利特检验}
    \begin{tblr}{
      hline{1,Z} ={solid},
      cell{1}{1} = {c=2}{c},
      cell{2}{1} = {r=3}{m}
    }
      KMO取样适切性量数 &      & 0.896    \\
      巴特利特球形度检验  & 近似卡方 & 2967.597 \\
                & 自由度  & 276      \\
                & 显著性  & .000   
    \end{tblr}
  \end{table}

  表~\ref{table:问卷效度分析：KMO和巴特利特检验} 显示，这份调查表的 KMO数值为0.896，高于0.7，适宜因素分析.另外，巴特利特测试值0.000适用于因素分析.
\end{enumerate}


\subsection{信度分析}

信度分析主要考察某一类问题的各项指标，其内在相似性和内在关联性，其高的信度表示所评定的某一类条目的可靠性高性.

对24个问题的调查结果进行了信度分析，结果显示在表~\ref{table:信度分析：可靠性统计}中。


\begin{table}[htbp]
  \centering
  \caption{信度分析：可靠性统计}
  \label{table:信度分析：可靠性统计}
  \begin{tblr}{
    hline{1,Z} ={solid},
    cells = {c}
  }
  克隆巴赫Alpha	& 基于标准化项的克隆巴赫Alpha &	项数 \\
  0.885	& 0.892	& 24 
  \end{tblr}
\end{table}


克朗巴赫指数0.885，高于0.7，说明该调查的可信性很高。



\section{调查问卷的结果分析}

本研究依据《课标2017》所订的《数学计算能力等级》的2-2架构架构，在每个等级的四个维度上，各有三个问题，共有24个问题，皆采用李氏五点式的测验，1代表非常赞同、2代表赞同、3代表不肯定、4代表反对、5代表极端反对.
在相同级别相同的情况下，三个问题的真实得分和平均值，平均得分显示了该研究人员在某个特定领域的工作情况.如果该领域的平均得分超过3，那么该领域的得分就会更好；如果平均得分在3以下，就会被视为是较坏的，而平均得分较低，说明该级别的这一部分的绩效较差.

初步调查和分析调查的执行状况：

性别：202名男生，187名女生.

班级分布：高中二年级随机抽取12个班，分别为：1班、2班、4班、5班、7班、8班、9班、10班、11班、12班、15班。

笔者将对389份有效问卷的结果进一步分析.


\subsection{高中生数学运算水平一分析}

\begin{enumerate}[(1)]
  \item 水平一情境与问题 \par
  测量数学运算水平一情境与问题的分别是第1、8、10题，具体统计结果见表~\ref{table:水平一情境与问题情况统计} 和图~\ref{figure:水平一情境与问题情况统计图}.
  \begin{table}[htbp]
    \centering
    \caption{水平一情境与问题情况统计}
    \label{table:水平一情境与问题情况统计}
    \begin{tblr}{
      width = \textwidth,
      hline{1,Z} ={solid},
      colspec = {X[11em,l]X[4em,c]X[2.5em,c]X[3em,c]X[3em,c]X[5em,c]X[3em,c]},
      cell{1}{1} = {c},
      cell{2}{1} = {r=2}{m},
      cell{2}{7} = {r=2}{m},
      cell{4}{1} = {r=2}{m},
      cell{4}{7} = {r=2}{m},
      cell{6}{1} = {r=2}{m},
      cell{6}{7} = {r=2}{m},
    }
      项目           & 非常同意    & 同意      & 不确定     & 不同意    & 非常不同意  & 平均分  \\
      1.我能够根据问题中的等量关系写出数学表达式 & 195     & 165     & 25      & 2      & 2      & 4.41 \\
          & 50.13\% & 42.42\% & 6.43\%  & 0.51\% & 0.51\% &      \\
      8.我在拿到数学题目后会认真审题 & 93      & 184     & 82      & 16     & 14     & 3.84 \\
            & 23.91\% & 47.30\% & 21.08\% & 4.11\% & 3.60\% &      \\
      {10.我能指岀某一数学问题是向量运算问题} & 102     & 125     & 124     & 26     & 12     & 3.72 \\
        & 26.22\% & 32.13\% & 31.88\% & 6.68\% & 3.08\% &     
    \end{tblr}
  \end{table}
  \begin{figure}[htbp]
    \centering
    \includegraphics[width = \textwidth]{3-1}
    \caption{水平一情境与问题情况统计图}
    \label{figure:水平一情境与问题情况统计图}
  \end{figure}

  从表~\ref{table:水平一情境与问题情况统计} 可以看出，第一道题的平均分数为4.41，第八道和第十道题的平均分数都在4左右，这表明苏中学的分数很高，这是因为这所学校的学生都已经具备了一定的数学能力……

  50.13\%的同学认为"我可以按照问题中的等值关系来书写算术公式"，而42.42\%的同学则表示赞同.

  至于“我会在收到数学问题后，仔细审查问题”，47.30\%的人赞同，23.91\%的人赞同，3.60\%的人不这么认为，所以，在认真的考题方面，高中阶段的同学们还有很大的不足，要养成一个好的作题习惯，必须要进-步培养.

  针对“我可以指一个…的数理问题是矢量法问题”，从资料中我们可以看到，赞成同意的人很多，32.13\%，不肯定的占31.88\%，反对的占3.08\%，认为这是一个很困难的问题。
  \item 水平一知识与技能 \par
  测量数学运算水平一知识与技能的分别是第3、11、16题，具体统计结果见表~\ref{table:水平一知识与技能情况统计} 和图~\ref{figure:水平一知识与技能情况统计图}.
  \begin{table}[htbp]
    \centering
    \caption{水平一知识与技能情况统计}
    \label{table:水平一知识与技能情况统计}
    \begin{tblr}{
      width = \textwidth,
      hline{1,Z} ={solid},
      colspec = {X[11em,l]X[4em,c]X[2.5em,c]X[3em,c]X[3em,c]X[5em,c]X[3em,c]},
      cell{1}{1} = {c},
      cell{2}{1} = {r=2}{m},
      cell{2}{7} = {r=2}{m},
      cell{4}{1} = {r=2}{m},
      cell{4}{7} = {r=2}{m},
      cell{6}{1} = {r=2}{m},
      cell{6}{7} = {r=2}{m},
    }
    项目                    & 非常同意    & 同意      & 不确定     & 不同意    & 非常不同意  & 平均分  \\
    3. 我能清楚记忆常见的数学公式和法则   & 144     & 187     & 50      & 6      & 2      & 4.2  \\
                          & 37.02\% & 48.07\% & 12.85\% & 1.54\% & 0.51\% &      \\
    11. 我能够选择相应的运算法则来进行计算 & 113     & 210     & 58      & 7      & 1      & 4.1  \\
                          & 29.05\% & 53.98\% & 14.91\% & 1.80\% & 0.26\% &      \\
    16. 在解题时，我更倾向算法则来进行计算 & 121     & 159     & 74      & 27     & 8      & 3.92 \\
                          & 31.11\% & 40.87\% & 19.02\% & 6.94\% & 2.06\% &     
    \end{tblr}
  \end{table}
  \begin{figure}[htbp]
    \centering
    \includegraphics[width = \textwidth]{3-2}
    \caption{水平一知识与技能情况统计图}
    \label{figure:水平一知识与技能情况统计图}
  \end{figure}

  表~\ref{table:水平一知识与技能情况统计} 表明，16号问题的平均得分是3.92，而第三和11的得分则稍微超过4。

  37.02\%的同学赞同“我能清晰地记住常用的数学公式和定律”，48.07\%的同学认为这一点是完全正确的.
  
  尽管“我可以根据自己的计算规则来进行计算”，但53.98\%的人赞同，29.05\%的人表示赞同，所以在计算规则的选取上，仍然存在着改进的余地。

  在回答问题时，我更喜欢考虑一段时间而不是马上开始写作，31.11\%的同学对这一问题持肯定态度，。在这种情况下，在学习过程中，学生可以按照问题的特点，形成相应的计算方法.
  
  \item 水平一思维与表达 \par
  测量数学运算水平一思维与表达的分别是第4、6、9题，具体统计结果见表~\ref{table:水平一思维与表达情况统计} 和图~\ref{figure:水平一思维与表达情况统计图}.
  \begin{table}[htbp]
    \centering
    \caption{水平一思维与表达情况统计}
    \label{table:水平一思维与表达情况统计}
    \begin{tblr}{
      width = \textwidth,
      hline{1,Z} ={solid},
      colspec = {X[10em,l]X[4em,c]X[2.5em,c]X[3em,c]X[3em,c]X[5em,c]X[3em,c]},
      cell{1}{1} = {c},
      cell{2}{1} = {r=2}{m},
      cell{2}{7} = {r=2}{m},
      cell{4}{1} = {r=2}{m},
      cell{4}{7} = {r=2}{m},
      cell{6}{1} = {r=2}{m},
      cell{6}{7} = {r=2}{m},
    }
    项目                         & 非常同意    & 同意      & 不确定     & 不同意    & 非常不同意  & 平均分  \\
    4. 在学习运算法则时，我 会思考该法则的意义和作用 & 93      & 147     & 115     & 29     & 5      & 3.76 \\
                                & 23.91\% & 37.79\% & 29.56\% & 7.46\% & 1.29\% &      \\
      6. 我能够运用运算验证简单的数学结论        & 143     & 177     & 52      & 14     & 3      & 4.4  \\
                                & 36.76\% & 45.50\% & 13.37\% & 3.60\% & 0.77\% &      \\
      9. 我能规范地书写运算过程             & 89      & 164     & 106     & 24     & 6      & 3.79 \\
                                & 22.88\% & 42.16\% & 27.25\% & 6.17\% & 1.54\% &     
    \end{tblr}
  \end{table}
  \begin{figure}[htbp]
    \centering
    \includegraphics[width = \textwidth]{3-3}
    \caption{水平一思维与表达情况统计图}
    \label{figure:水平一思维与表达情况统计图}
  \end{figure}

  从表~\ref{table:水平一思维与表达情况统计} 可以看到，这三道题目的平均分数都在4左右，而第四和第九道题目的平均值低于4，就数学的计算能力而言，从知识和技巧上来讲，数学的计算能力，思考和表达能力都很弱，但高中生的思维与表达能力仍需进一步提升。


  29.56\%的人不太肯定"在计算定律中我会考虑其重要性和功能"，而不赞同的比例为7.46\%.

  在“我可以使用计算来证明一个简单的数学结果”时，82\%的同学持赞同或更高的观点，只有0.77\%的人不这么认为，大多数人都可以通过计算来证明一个简单的数学结果。

  对于“我能够正确地写出计算程序”的学生中，42.16\%的人认为“我能正确地写出计算程序”，但是27.25\%的人不肯定，6.17\%的人不这么认为.所以，在标准的书写计算这一运算习惯方面表现并不出色，规范书写运算过程会影响最后得分，需要引起学生和教师的重视.

  \item 横向一沟通和思考 \par
  计算数学操作等级1的第二题、14题和18题分别为沟通和思考，详细数据如表~\ref{table:横向一沟通和思考情况统计} 和图~\ref{figure:横向一沟通和思考情况统计图} 所示.
  \begin{table}[htbp]
    \centering
    \caption{横向一沟通和思考情况统计}
    \label{table:横向一沟通和思考情况统计}
    \begin{tblr}{
      width = \textwidth,
      hline{1,Z} ={solid},
      colspec = {X[11em,l]X[4em,c]X[2.5em,c]X[3em,c]X[3em,c]X[5em,c]X[3em,c]},
      cell{1}{1} = {c},
      cell{2}{1} = {r=2}{m},
      cell{2}{7} = {r=2}{m},
      cell{4}{1} = {r=2}{m},
      cell{4}{7} = {r=2}{m},
      cell{6}{1} = {r=2}{m},
      cell{6}{7} = {r=2}{m},
    }
    项目                          & 非常同意    & 同意      & 不确定     & 不同意     & 非常不同意  & 平均分  \\
    2. 我能对熟悉的数学应用与建模的结果作出说明 或解释 & 100     & 167     & 102     & 17      & 3      &   3.88   \\
                                & 25.71\% & 42.93\% & 26.22\% & 4.37\%  & 0.77\% &  \\
    14. 我平时在做完数学运算后，会再检查一下      & 81      & 143     & 106     & 42      & 17     & 3.59 \\
                                & 20.82\% & 36.76\% & 27.25\% & 10.80\% & 4.37\% &      \\
    18. 我会与同学交流某道题的运算结果         & 133     & 178     & 50      & 21      & 7      & 4.05 \\
                                & 34.19\% & 45.76\% & 12.85\% & 5.40\%  & 1.80\% &    
    \end{tblr}
  \end{table}
  \begin{figure}[htbp]
    \centering
    \includegraphics[width = \textwidth]{3-4}
    \caption{水平一交流与反思情况统计图}
    \label{figure:水平一交流与反思情况统计图}
  \end{figure}

  如表~\ref{table:横向一沟通和思考情况统计} 所示，第2题的平均得分为3.88,18题为4.05，两者皆为交流，得分较高，14题为3.59，为上述四项中最高.

  对于“我能对数理的运用及模式做出的解释”，认同率为42.93\%，持怀疑态度的占26.22\%，赞成“与我的好友交流一道问题的答案”的比例达到45.19\%.

  至于“我一般会在完成一道算术后，会重新核对一遍”，百分之十五以上的人都不这么认为，所以，学生们在检查和反省上仍有不足之处，他们的计算能力并没有被培养出来，而通过实际的研究，他们的计算习惯也是如此.
\end{enumerate}


\subsection{高中生数学运算水平二分析}

\begin{enumerate}[(1)]
  \item 水平二情境与问题 \par
  测量数学运算水平二情境与问题的分别是第12、20、22题，具体统计结果见表~\ref{table:水平二情境与问题情况统计} 和图~\ref{figure:水平二情境与问题情况统计图}.
  \begin{table}[htbp]
    \centering
    \caption{水平二情境与问题情况统计}
    \label{table:水平二情境与问题情况统计}
    \begin{tblr}{
      width = \textwidth,
      hline{1,Z} ={solid},
      colspec = {X[11em,l]X[4em,c]X[2.5em,c]X[3em,c]X[3em,c]X[5em,c]X[3em,c]},
      cell{1}{1} = {c},
      cell{2}{1} = {r=2}{m},
      cell{2}{7} = {r=2}{m},
      cell{4}{1} = {r=2}{m},
      cell{4}{7} = {r=2}{m},
      cell{6}{1} = {r=2}{m},
      cell{6}{7} = {r=2}{m},
    }
    项目                         & 非常同意    & 同意      & 不确定     & 不同意    & 非常不同意  & 平均分  \\
    12. 我能够将运算方法运用到相关情境中解决问题   & 80      & 180     & 107     & 20     & 2      & 3.81 \\
                                & 20.57\% & 46.27\% & 27.51\% & 5.14\% & 0.51\% &      \\
      20. 在解决有关运算问题时，我能联想类似的运算方法 & 68      & 142     & 154     & 22     & 3      & 3.64 \\
                                & 17.48\% & 36.50\% & 39.59\% & 5.66\% & 0.77\% &      \\
      22. 我会自己尝试总结一类问题的运算技巧或方法   & 83      & 145     & 121     & 36     & 4      & 3.69 \\
                                & 21.34\% & 37.28\% & 31.11\% & 9.25\% & 1.03\% &  
    \end{tblr}
  \end{table}
  \begin{figure}[htbp]
    \centering
    \includegraphics[width = \textwidth]{3-5}
    \caption{水平二情境与问题情况统计图}
    \label{figure:水平二情境与问题情况统计图}
  \end{figure}

  表~\ref{table:水平二情境与问题情况统计} 表明，与数学操作等级一的情景与问题比较，第二个问题的得分要低一些，20和22的分数都在3.5.

  46.27\%的同学说：“我可以在有关情况下使用计算法来解问题”，持相同意的同学有46.27\%，成绩还可以.

  在处理相关操作问题时，只有17.48\%的同学表示赞同，39.59\%的同学认为不一定，说明在处理操作问题时，计算法的联系不强，但整体的水准还是可以接受的.

  至于「我会自行归纳一种计算的技术或方式」，接近60\%的同学表示赞同及以上，但不肯定的占31.11\%,9.25\%的同学则不认同，在此题目中，反对的情绪更为强烈，显示出他们在归纳某一类别问题时表现不佳.

  整体来说，在相关的环境下，学生决定操作目标，并提出操作问题，60\%以上的同学已经基本满足了高考的目标，但是，怎样从水平、水平二跨越，值得一线教师做进一步探讨.

  \item 水平二知识与技能 \par
  测量数学运算水平二知识与技能的分别是第7、15、24题，具体统计结果见表~\ref{table:水平二知识与技能情况统计} 和图~\ref{figure:水平二知识与技能情况统计图}.
  \begin{table}[htbp]
    \centering
    \caption{水平二知识与技能情况统计}
    \label{table:水平二知识与技能情况统计}
    \begin{tblr}{
      width = \textwidth,
      hline{1,Z} ={solid},
      colspec = {X[11em,l]X[4em,c]X[2.5em,c]X[3em,c]X[3em,c]X[5em,c]X[3em,c]},
      cell{1}{1} = {c},
      cell{2}{1} = {r=2}{m},
      cell{2}{7} = {r=2}{m},
      cell{4}{1} = {r=2}{m},
      cell{4}{7} = {r=2}{m},
      cell{6}{1} = {r=2}{m},
      cell{6}{7} = {r=2}{m},
    }
    项目                         & 非常同意    & 同意      & 不确定     & 不同意    & 非常不同意  & 平均分  \\
    7. 在进行数学运算时,我能合理选择公式，并灵活运用      & 91      & 172     & 106     & 18     & 2      & 3.85 \\
    & 23.39\% & 44.22\% & 27.25\% & 4.63\% & 0.51\% &      \\
    15. 当运算进行不下去时，我会改变算法或者寻求别的方法    & 145     & 187     & 42      & 12     & 3      &      \\
        & 37.28\% & 4\&07\% & 10.80\% & 3.08\% & 0.77\% &      \\
    24. 我能够针对不同类型问题的特点，选择合理简洁的 运算途径 & 65      & 142     & 150     & 26     & 6      & 3.60 \\
        & 16.71\% & 36.50\% & 38.56\% & 6.68\% & 1.54\% &  
    \end{tblr}
  \end{table}
  \begin{figure}[htbp]
    \centering
    \includegraphics[width = \textwidth]{3-6}
    \caption{水平二知识与技能情况统计图}
    \label{figure:水平二知识与技能情况统计图}
  \end{figure}

  表~\ref{table:水平二知识与技能情况统计} 表明，在数学操作能力2中，在知识和技巧上的成绩最好，与其在数学操作能力上的成绩是一样的.

  关于“在做算术时，我能正确地选取公式，并能灵活应用”，赞成的比例为44.22\%，但不肯定的占27.25\%，所以，计算公式的弹性方面基本上达到了要求，但还需要进一步改.

  有85\%的同学认为，在计算不出结果时，可以修改算法，也可以寻找其他办法.

  在“我可以根据各种问题的特征，选取最简单、最简单的计算方式”的情况下，只有16.71\%的人表示完全赞同，38.56\%的人认为，在计算精练和灵活性上，还有很大的差距。

  总的来说，他们的学习能力和能力都很强，可以根据自己想要的问题，选择合理的、简洁的计算方式.
  \item 水平二思维与表达 \par
  测量数学运算水平二思维与表达的分别是第17、19、23题，具体统计结果见表~\ref{table:水平二思维与表达情况统计} 和图~\ref{figure:水平二思维与表达情况统计图}.
  \begin{table}[htbp]
    \centering
    \caption{水平二思维与表达情况统计}
    \label{table:水平二思维与表达情况统计}
    \begin{tblr}{
      width = \textwidth,
      hline{1,Z} ={solid},
      colspec = {X[11em,l]X[4em,c]X[2.5em,c]X[3em,c]X[3em,c]X[5em,c]X[3em,c]},
      cell{1}{1} = {c},
      cell{2}{1} = {r=2}{m},
      cell{2}{7} = {r=2}{m},
      cell{4}{1} = {r=2}{m},
      cell{4}{7} = {r=2}{m},
      cell{6}{1} = {r=2}{m},
      cell{6}{7} = {r=2}{m},
    }
    项目                       & 非常同意    & 同意      & 不确定     & 不同意     & 非常不同意  & 平均分  \\
    17. 我能够利用分类讨论解决复杂的代数运算问题 & 82      & 205     & 90      & 11      & 1      & 3.92 \\
                              & 21.08\% & 52.70\% & 23.14\% & 2.83\%  & 0.26\% &      \\
      19. 在综合运算问题中，我会因运算不规范而失分 & 74      & 118     & 98      & 77      & 22     & 2.63 \\
                              & 19.02\% & 30.33\% & 25.19\% & 19.79\% & 5.66\% &      \\
      23. 我能由数思形思考代数运算问题       & 83      & 153     & 127     & 25      & 1      & 3.75 \\
                              & 21.34\% & 39.33\% & 32.65\% & 6.43\%  & 0.26\% &   
    \end{tblr}
  \end{table}
  \begin{figure}[htbp]
    \centering
    \includegraphics[width = \textwidth]{3-7}
    \caption{水平二思维与表达情况统计图}
    \label{figure:水平二思维与表达情况统计图}
  \end{figure}

  表~\ref{table:水平二思维与表达情况统计} 表明第19题的平均分数只有2.63 （为符合其它学生的成绩，都是积极的，在进行运算时，十分赞同1，赞同2，如此等等），第17题和23题的平均值都是4分.

  针对「在「计算错误造成的」「计算错误」，有19.79\%的同学表达了不同意，5.66\%的同学则表示强烈反对，故在「全面问题」中，作业标准问题表现得尤为明显，并进一步显示了「计算习性」问题，应该予以重视.

  至于「我可以运用分类探讨来解复杂的代数问题」，52.70\%的同学持同意，21.08\%的同学则表示赞同，并取得很好的成绩.
  其中39.33\%的同学赞同“用数学思维来解决代数学问题”，32.65\%的同学认为这是个问题.

  通常情况下，学生在处理复杂的数学问题时，如果不能充分运用迁移、联想等思维方式，就会遇到一些复杂的计算问题。


  \item 水平二交流与反思 \par
  测量数学运算水平二交流与反思的分别是第5、13、21题，具体统计结果见表~\ref{table:水平二交流与反思情况统计} 和图~\ref{figure:水平二交流与反思情况统计图}.
  \begin{table}[htbp]
    \centering
    \caption{水平二交流与反思情况统计}
    \label{table:水平二交流与反思情况统计}
    \begin{tblr}{
      width = \textwidth,
      hline{1,Z} ={solid},
      colspec = {X[11em,l]X[4em,c]X[2.5em,c]X[3em,c]X[3em,c]X[5em,c]X[3em,c]},
      cell{1}{1} = {c},
      cell{2}{1} = {r=2}{m},
      cell{2}{7} = {r=2}{m},
      cell{4}{1} = {r=2}{m},
      cell{4}{7} = {r=2}{m},
      cell{6}{1} = {r=2}{m},
      cell{6}{7} = {r=2}{m},
    }
    项目                       & 非常同意    & 同意      & 不确定     & 不同意     & 非常不同意  & 平均分  \\
    5. 在完成数学运算后，我会思考其他的运算方法      & 39      & 124     & 142     & 66      & 18     & 3.26 \\
    & 10.03\% & 31.88\% & 36.50\% & 16.97\% & 4.63\% &      \\
    13. 我能够和同学交流思路，探讨更为合理简洁的运算方法 & 85      & 160     & 101     & 35      & 8      &  3.95    \\
        & 21.85\% & 41.13\% & 25.96\% & 9.00\%  & 2.06\% & 3.72 \\
    21. 对运算中的错误，我会寻找原因，并进行反思     & 99      & 203     & 74      & 10      & 3      & 3.98 \\
        & 25.45\% & 52.19\% & 19.02\% & 2.57\%  & 0.77\% &  
    \end{tblr}
  \end{table}
  \begin{figure}[htbp]
    \centering
    \includegraphics[width = \textwidth]{3-8}
    \caption{水平二交流与反思情况统计图}
    \label{figure:水平二交流与反思情况统计图}
  \end{figure}

  从表~\ref{table:水平二交流与反思情况统计} 可知，计算等级2的沟通与思考能力依然很弱，与在计算能力1中的成绩大致相同，第五个问题的平均值只有3.26,13和21的得分也是如此.

  针对「计算中的失误，我会找出其中的缘由，并加以检讨」，接近80\%的同学认同及上述观点，而学生则比较满足地找到问题的成因.
  41.13\%的学生赞同“我可以和同学们进行思想沟通，讨论更加科学、简明的计算方式”，他们的成绩还可以，但相对于水平中学的沟通来说，他们还是比较薄弱的.

  至于“做完后，我会考虑其它计算方式”，只有10.03\%的人赞同，绝大多数人持中立观点，说明很少有人会在做完一道题后，就开始考虑其它的问题，从而影响到他们的计算能力.

  总之，在沟通上还需要提高，在沟通中要注意利用计算进行讨论.
  

\end{enumerate}


\subsection{高中生数学运算能力上的表现及性别差异分析}

本研究共收集到389份有效的问卷，经统计，男性202名，女生187名。从总体和性别两个层面，对高中生数学计算能力的不同进行了研究。


在同一层次、同一维度设置3道题，由考生综合成绩，平均平均分代表考生在某一层次上的表现。男生、女生和总体 之间的各维度表现差异如下表~\ref{table:高中男女生数学运算能力表现的差异分析} 所示.

\begin{table}[htbp]
  \centering
  \caption{高中男女生数学运算能力表现的差异分析}
  \label{table:高中男女生数学运算能力表现的差异分析}
  \begin{tblr}{
    % hlines,vlines,
    hline{1,Z} = {solid},
    cells = {c},
    cell{2}{1,5,7,8} = {r=2}{m},
    cell{4}{1,5,7,8} = {r=2}{m},
    cell{6}{1,5,7,8} = {r=2}{m},
    cell{8}{1,5,7,8} = {r=2}{m},
    cell{10}{1,5,7,8} = {r=2}{m},
    cell{12}{1,5,7,8} = {r=2}{m},
    cell{14}{1,5,7,8} = {r=2}{m},
    cell{16}{1,5,7,8} = {r=2}{m},
    cell{1}{4} = {c=2}{c}
  }
      方面 &性别   & N     &  均值&	 &  标准差       &  T    &  P\\
    情境与问题 & 男 & 202 & 4.0363 & 3.989 & 0.68692  & 1.502 & 0.134 \\
      & 女 & 187 & 3.9376 &       & 0.60189 &       &       \\
    知识与技能 & 男 & 202 & 4.1436 & 4.071 & 0.63242 & 2.513 & 0.012 \\
          & 女 & 187 & 3.9929 &       & 0.54974 &       &       \\
    思维与表达 & 男 & 202 & 4.0264 & 3.894 & 0.60161 & 4.55  & .000     \\
          & 女 & 187 & 3.7504 &       & 0.59327 &       &       \\
    交流与反思 & 男 & 202 & 3.962  &   3.841    & 0.66848 & 4.085 & .000     \\
          & 女 & 187 & 3.7112 &       & 0.53973 &       &       \\
    情境与问题 & 男 & 202 & 3.8779 & 3.714 & 0.69544 & 5.233 & .000     \\
          & 女 & 187 & 3.5365 &       & 0.58989 &       &       \\
    知识与技能 & 男 & 202 & 3.9884 & 3.878 & 0.72267 & 3.508 & 0.001 \\
          & 女 & 187 & 3.7594 &       & 0.56031 &       &       \\
    思维与表达 & 男 & 202 & 3.5314 & 3.431 & 0.66613 & 3.392 & 0.001 \\
          & 女 & 187 & 3.3226 &       & 0.57932 &       &       \\
    交流与反思 & 男 & 202 & 37,624 & 3.655 & 0.6844  & 3.392 & 0.001 \\
          & 女 & 187 & 3.5383 &       & 0.61263 &       &   
  \end{tblr}
\end{table}

从表~\ref{table:高中男女生数学运算能力表现的差异分析} 可以看到，在各个层次，学生的知识和技巧平均为4.071，在各个层次中，他们的知识和技巧平均为3.878，高于其他三个层面.

结果显示，男生的数学计算成绩比女生要好，这是因为男女学生的数学记忆力和思考方式不同，所以，男孩和女孩的数学计算成绩比女生要好，这是因为男女学生的数学记忆力和思考方式的不同.

我们还进行了一次方差分析，结果表明： F统计是39.756，相应的测试有统计学意义 P<0.05，这表明了各维度分数有显著性的差别. LSD结果表明：水平一的环境与问题、知识和技巧的分数没有显著差别（均值为负且 P>0.05)，而且与其他两个方面（均值为正且 P<0.05). 水平一的情景与沟通、思考、思考、沟通、思考等方面的得分均未有统计学意义（均值为负且 P>0.05). 在情境和问题、沟通和反思上，没有明显的差异（均值为正， P>0.05）。

从整体上看，大学生在学习和学习能力上的成绩较好，但在思考能力和表达能力上则较弱，随着能力的提高，情况与问题、知识和技巧的差距也越来越大，而在第二阶段，他们的学习能力和能力的差距也越来越大.

\begin{table}[htbp]
  \centering
  \caption{不同维度单因素方差分析}
  \label{table:不同维度单因素方差分析}
  \begin{tblr}{
    hline{1,2,Z} = {solid},
    cells = {c}
  }
    & 平方和      & df   & 均方     & F      & 显著性 \\
    组间 & 113.041  & 7    & 16.149 & 39.756 & .000   \\
    组内 & 1264.091 & 3112 & 0.406  &        &     \\
    总数 & 1377.133 & 3119 &        &        &    
  \end{tblr}
\end{table}
% 进一步进行单因素方差分析，由表~\ref{table:不同维度单因素方差分析} 可知，F统计量为39.756,对应的检验 显著性P为0.000 （<0.05）,说明不同维度得分之间存在显著差异.我们进一步通过 LSD事后比较发现，水平一情境与问题和知识与技能的得分无显著差异（均值为正 且P>0.05）,并且明显高于水平一的另外两个方面（均值为正且P<0.05）,而水平一 思维与表达和交流与反思方面也无显著差异（均值为正且P>0.05）.水平二知识与技 能方面的得分明显高于水平二的另外三个方面（均值为正且P<0.05）.水平二的情境 与问题和交流与反思方面的得分无显著差异（均值为正且P>0.05）,并且明显高于思 维与表达方面（均值为正且P<0.05）.

% 综上可知，学生在知识与技能方面中表现突出，在思维与表达方面表现不佳.随着水平的上升，学生在情境与问题和知识与技能两方面表现的差异逐渐显现，由原来 的无显著差异，转变为水平二中知识与技能方面表现更为突出.同样，在水平一中无 显著差异的思维与表达和交流与反思，在水平二中的差异明显，思维与表达更显落 后.因此，笔者猜想水平的递进可能会拉开各方面表现的距离.



\section{对试卷的错题分析}

由图~\ref{figure:4-1} 中的第2题可以发现学生没有看清楚题目要求，没有好好审题,也有可能是学生没有记清楚使用基本不等式的条件、乱用公式。教师在讲解知识点时，要注意运用不等式的条件，要记住"一正二定三相等"。


\begin{figure}[htbp]
  \centering
  \includegraphics[width = 0.7\textwidth]{4-1}
  \caption{}
  \label{figure:4-1}
\end{figure}

由图~\ref{figure:4-2} 中的第8题可以看出填空的答案是要写成集合的形式，虽然学生算出来的答案是对的，但是不注意书写，最终的答案还是错的。学生平时做题时书写一定要规范，一定要知道集合的书写格式。

\begin{figure}[htbp]
  \centering
  \includegraphics[width = 0.7\textwidth]{4-2}
  \caption{}
  \label{figure:4-2}
\end{figure}

由图~\ref{figure:4-3} 可以看出这是学生经常会出错的问题，粗心大意、抄错内容，这种错误对于高中生来说太不应该了。平时做题时候一定要注意千万不能抄错，考试中一定要避免因为这种小错误而丢分现象。

\begin{figure}[htbp]
  \centering
  \includegraphics[width = 0.7\textwidth]{4-3}
  \caption{}
  \label{figure:4-3}
\end{figure}



\newpage
\section{访谈内容}

笔者对西藏拉萨中学的几位老师进行了访谈来更好地了解西藏高中生数学运算能力的现状，以下是两位老师的访谈反馈：

\begin{enumerate}[\chinese*、]
  \item 根据您多年的教学经验，您觉得现在的高中生在数学学习中存在哪些问题？
    \begin{enumerate}[A:]
      \item 概念的理解和抽象提炼，数学运算，分析解决问题，数学学习习惯和能力。
      \item 我认为目前西藏高中生在学习数学的过程中主要存在问题
        \begin{enumerate}[1 ,leftmargin = 1em]
          \item 对教材的解读和掌握不透彻;
          \item 对公式，概念，性质的理解及应用不充分，不会学以致用;
          \item 在实际的做题过程中做题能力及计算能力差;
          \item 还有多种原因使得藏区高中生的数学基础非常薄弱
        \end{enumerate}
    \end{enumerate}
  \item 您认为现阶段西藏高中生的数学运算能力存在哪些问题?
    \begin{enumerate}[A:]
      \item 数学运算普遍较差，大部分学生不能很灵活地运用各种公式及没有良好的运算习惯。
      \item 在数学运算当中运算能力差，对一些公式的应用不够灵活。 在数系的扩充当中唯一不变的是运算性质，但在学生心里学了新的内容就觉得前面的一些运算性质，公式是不能用的。 比如在必修一当中涉及到函数的单调性可以用“五步法”来判断，但学完导数后学生认为判断函数的单调性只能用导数，但事实是前面的方法依然可以用，总的来说缺乏贯通、灵活的思想。
    \end{enumerate}
  \item 您认为教师日常教学中重视学生运算能力的培养吗？
    \begin{enumerate}[A:]
      \item 大多数教师都比较很注重运算和书写规范。
      \item 因人而异，有些老师非常注重，但有些老师不管，认为运算是学生自身的问题，应该靠自己提升，但实际上每一个数学老师都有提高学生运算能力的责任，尤其在西藏任教的老师。
    \end{enumerate}
  \item 您认为应该怎样有效地培养西藏高中生数学运算能力？
    \begin{enumerate}[A:]
      \item 经常性创设适当的限时限量答题训练，提高学生专注度和答题效率。长期训练有一定改善。
      \item 我认为提高运算能力首先要掌握涉及高中内容的一些公式，其次是要做题，见多识广，题目做多了，也就掌握考法，运算能力也就上去了。
    \end{enumerate}
\end{enumerate}
